(1) 界面切向波数守恒，\(k_x=n_1k_0\sin\theta\)。入射侧法向分量\(k_{z1}=\sqrt{n_1^2k_0^2-k_x^2}=n_1k_0\cos\theta>0\)，透射侧在全反射区为\(k_{z2}=i\kappa\)，\(\kappa=\sqrt{k_x^2-n_2^2k_0^2}>0\)；因此低折射率侧场沿法线指数衰减而无传播法向功率。
(2) 对TE偏振，Fresnel反射系数\(r_s=(k_{z1}-k_{z2})/(k_{z1}+k_{z2})=(k_{z1}-i\kappa)/(k_{z1}+i\kappa)\)，故\(|r_s|=1\)，相位可取\(\phi_r(k_x)=-2\arctan(\kappa/k_{z1})\)。虽反射率为1，但相位随入射波数变化，形成反射波包空间移动的机制。
(3) 将入射束作围绕中心\(k_x\)的角谱积分，反射振幅额外含\(e^{i\phi_r(k_x)}\)。对窄角谱作驻相，光斑位移\(D=-d\phi_r/dk_x\)。利用\(dk_{z1}/dk_x=-k_x/k_{z1}\)、\(d\kappa/dk_x=k_x/\kappa\)，得到\(d[\arctan(\kappa/k_{z1})]/dk_x=k_x/(\kappa k_{z1})\)，故\(D=2k_x/(\kappa k_{z1})\)，单位是长度且为正。
(4) 当\(\theta\downarrow\theta_c\)，\(\kappa\to0\)使平面波公式的\(D\)形式发散；当\(\theta\to\pi/2\)，\(k_{z1}\to0\)也会产生形式发散。真实有限宽度束在临界角附近有宽角谱，不能将反射相位仅作一阶Taylor展开，需用完整角谱积分；束宽、偏振和材料吸收都会使位移有限或改变峰位。
